\documentclass[a4paper, 12pt]{article}

\usepackage[utf8]{inputenc}
\usepackage[T2A]{fontenc}
\usepackage[english,russian]{babel}
\usepackage{amsmath, amsfonts, amsthm, amssymb, mathtools}
\usepackage{multicol}
\usepackage{enumitem}
\usepackage{cancel}
\usepackage{hyperref}
\hypersetup{
    colorlinks,
    linkcolor={black},
    citecolor={black},
    urlcolor={blue},
}
\usepackage{indentfirst}
\usepackage{adjustbox}
\usepackage{tikz}
\usepackage{float}
\usepackage[table]{xcolor}

\newcommand{\re}{\mathbf{r}}

\title{Лабораторная работа №2}
\author{Вариант 26}

\begin{document}
    \maketitle
    \tableofcontents
    \newpage

    \section{Задание}

    Дано регулярное выражение $\re$:

    \[
       (aba)^{*}((aa | bab)aba^{*})^{*}
    \]

    По предоставленному регулярному выражению построить:

    \begin{itemize}
        \itemsep -0.8pt
        \item[---] минимальный ДКА;
        \item[---] возможно малый НКА;
        \item[---] возможно малый переключающийся (с конъюнкцией) КА;
        \item[---] расширенное регулярное выражение.
    \end{itemize}

    \section{Недетерминированный конечный автомат}

    Для исходного регулярного выражения построим автомат Глушкова. Получаем НКА с 12 состояниями.

    %% Строим через Chipollino.
    %% Glushkov {(aba)*((aa|bab)aba*)*} !!

    \begin{figure}[H]
      \centering
      \begin{adjustbox}{width=\textwidth}
      
\def \globalscale {1.000000}
\begin{tikzpicture}[y=1cm, x=1cm, yscale=\globalscale,xscale=\globalscale, every node/.append style={scale=\globalscale}, inner sep=0pt, outer sep=0pt]
  \begin{scope}[shift={(0.1411, -9.8707)}]
    \path[fill=white] (-0.1411, 9.8778) -- (-0.1411, 19.8896) -- (38.4993, 19.8896) -- (38.4993, 9.8778) -- (-0.1411, 9.8778) -- (-0.1411, 9.8778)-- cycle;;



    \path[draw=black,fill=black] (0.0635, 16.3336) ellipse (0.0635cm and 0.0635cm);



    \path[draw=black] (2.2084, 16.3336) ellipse (0.635cm and 0.635cm);



    \path[draw=black] (2.2084, 16.3336) ellipse (0.7761cm and 0.7761cm);



    \node[anchor=south] (text4466) at (2.2084, 16.1854){S};



    \path[draw=black] (0.1348, 16.3336).. controls (0.2501, 16.3336) and (0.6265, 16.3336) .. (1.0252, 16.3336);



    \path[draw=black,fill=black] (1.016, 16.4571) -- (1.3688, 16.3336) -- (1.016, 16.2101) -- (1.016, 16.4571) -- (1.016, 16.4571)-- cycle;;



    \path[draw=black] (5.2698, 14.4639) ellipse (0.7962cm and 0.7962cm);



    \node[anchor=south] (text8733) at (5.2698, 14.3157){a.0};



    \path[draw=black] (2.8123, 15.839).. controls (3.0526, 15.6422) and (3.3412, 15.4217) .. (3.6195, 15.2471).. controls (3.7934, 15.138) and (3.9843, 15.0333) .. (4.1727, 14.938);



    \path[draw=black,fill=black] (4.2136, 15.0555) -- (4.4771, 14.7906) -- (4.1056, 14.8332) -- (4.2136, 15.0555) -- (4.2136, 15.0555)-- cycle;;



    \node[anchor=south] (text1305) at (3.7292, 15.3952){a};



    \path[draw=black] (34.4809, 13.5114) ellipse (0.7962cm and 0.7962cm);



    \node[anchor=south] (text9095) at (34.4809, 13.3632){a.3};



    \path[draw=black] (2.5852, 17.0169).. controls (3.0572, 17.852) and (3.9924, 19.1558) .. (5.2345, 19.1558).. controls (5.2345, 19.1558) and (5.2345, 19.1558) .. (31.1231, 19.1558).. controls (33.1728, 19.1558) and (33.9831, 16.3583) .. (34.2833, 14.702);



    \path[draw=black,fill=black] (34.4039, 14.7292) -- (34.3401, 14.3609) -- (34.1602, 14.6883) -- (34.4039, 14.7292) -- (34.4039, 14.7292)-- cycle;;



    \node[anchor=south] (text7206) at (18.0093, 19.304){a};



    \path[draw=black] (14.9087, 14.1464) ellipse (0.8153cm and 0.8153cm);



    \node[anchor=south] (text5117) at (14.9087, 13.9982){b.5};



    \path[draw=black] (2.9933, 16.3336).. controls (3.5983, 16.3336) and (4.4711, 16.3336) .. (5.2345, 16.3336).. controls (5.2345, 16.3336) and (5.2345, 16.3336) .. (11.6748, 16.3336).. controls (12.6287, 16.3336) and (13.5072, 15.6626) .. (14.1012, 15.064);



    \path[draw=black,fill=black] (14.181, 15.1596) -- (14.3327, 14.8181) -- (14.001, 14.9902) -- (14.181, 15.1596) -- (14.181, 15.1596)-- cycle;;



    \node[anchor=south] (text6491) at (8.3979, 16.4818){b};



    \path[draw=black] (8.3979, 12.3825) ellipse (0.8153cm and 0.8153cm);



    \node[anchor=south] (text5133) at (8.3979, 12.2343){b.1};



    \path[draw=black] (5.9044, 13.9714).. controls (6.1468, 13.7791) and (6.4322, 13.5611) .. (6.7007, 13.3773).. controls (6.9, 13.2412) and (7.118, 13.1029) .. (7.3286, 12.9748);



    \path[draw=black,fill=black] (7.3858, 13.0845) -- (7.626, 12.7981) -- (7.2595, 12.8722) -- (7.3858, 13.0845) -- (7.3858, 13.0845)-- cycle;;



    \node[anchor=south] (text8912) at (6.8241, 13.5255){b};



    \path[draw=black] (11.6395, 12.9117) ellipse (0.7962cm and 0.7962cm);



    \path[draw=black] (11.6395, 12.9117) ellipse (0.9373cm and 0.9373cm);



    \node[anchor=south] (text6285) at (11.6395, 12.7635){a.2};



    \path[draw=black] (9.217, 12.5134).. controls (9.5483, 12.5688) and (9.943, 12.6344) .. (10.3152, 12.6965);



    \path[draw=black,fill=black] (10.2852, 12.8168) -- (10.6535, 12.7529) -- (10.3258, 12.573) -- (10.2852, 12.8168) -- (10.2852, 12.8168)-- cycle;;



    \node[anchor=south] (text9494) at (9.9579, 12.7952){a};



    \path[draw=black] (10.7294, 13.1833).. controls (10.5135, 13.2475) and (10.2828, 13.3135) .. (10.0672, 13.3703).. controls (8.8399, 13.6941) and (7.4207, 14.0127) .. (6.4527, 14.2219);



    \path[draw=black,fill=black] (6.4389, 14.0988) -- (6.1196, 14.2935) -- (6.4908, 14.3401) -- (6.4389, 14.0988) -- (6.4389, 14.0988)-- cycle;;



    \node[anchor=south] (text4443) at (8.3979, 14.1118){a};



    \path[draw=black] (12.1221, 12.0904).. controls (12.6545, 11.2353) and (13.6306, 10.0189) .. (14.8735, 10.0189).. controls (14.8735, 10.0189) and (14.8735, 10.0189) .. (31.1231, 10.0189).. controls (32.4721, 10.0189) and (33.4398, 11.4039) .. (33.9711, 12.4294);



    \path[draw=black,fill=black] (33.8571, 12.4778) -- (34.1235, 12.7399) -- (34.079, 12.3687) -- (33.8571, 12.4778) -- (33.8571, 12.4778)-- cycle;;



    \node[anchor=south] (text7220) at (22.7111, 10.1671){a};



    \path[draw=black] (12.5321, 13.2426).. controls (12.9106, 13.3886) and (13.3608, 13.5625) .. (13.7643, 13.7181);



    \path[draw=black,fill=black] (13.7104, 13.8296) -- (14.0839, 13.8416) -- (13.7993, 13.5992) -- (13.7104, 13.8296) -- (13.7104, 13.8296)-- cycle;;



    \node[anchor=south] (text2852) at (13.335, 13.7142){b};



    \path[draw=black] (37.562, 12.2061) ellipse (0.7962cm and 0.7962cm);



    \node[anchor=south] (text61) at (37.562, 12.0579){a.4};



    \path[draw=black] (35.2294, 13.202).. controls (35.5932, 13.0443) and (36.0433, 12.8492) .. (36.4483, 12.6739);



    \path[draw=black,fill=black] (36.4949, 12.7882) -- (36.7693, 12.5345) -- (36.3968, 12.5614) -- (36.4949, 12.7882) -- (36.4949, 12.7882)-- cycle;;



    \node[anchor=south] (text5828) at (36.0214, 13.0461){a};



    \path[draw=black] (24.2658, 12.065) ellipse (0.7962cm and 0.7962cm);



    \node[anchor=south] (text4952) at (24.2658, 11.9168){a.8};



    \path[draw=black] (36.7792, 11.9863).. controls (36.1745, 11.8279) and (35.2979, 11.6417) .. (34.5161, 11.6417).. controls (27.4997, 11.6417) and (27.4997, 11.6417) .. (27.4997, 11.6417).. controls (26.8171, 11.6417) and (26.0593, 11.7376) .. (25.4529, 11.8382);



    \path[draw=black,fill=black] (25.4385, 11.7154) -- (25.1125, 11.8981) -- (25.4811, 11.9585) -- (25.4385, 11.7154) -- (25.4385, 11.7154)-- cycle;;



    \node[anchor=south] (text6284) at (31.0878, 11.7898){a};



    \path[draw=black] (18.0093, 12.2061) ellipse (0.7962cm and 0.7962cm);



    \node[anchor=south] (text7455) at (18.0093, 12.0579){a.6};



    \path[draw=black] (15.5614, 13.6363).. controls (15.8037, 13.4458) and (16.0881, 13.2338) .. (16.359, 13.0598).. controls (16.5379, 12.9452) and (16.734, 12.8319) .. (16.9256, 12.7272);



    \path[draw=black,fill=black] (16.982, 12.8372) -- (17.2367, 12.5635) -- (16.867, 12.6185) -- (16.982, 12.8372) -- (16.982, 12.8372)-- cycle;;



    \node[anchor=south] (text3209) at (16.4687, 13.208){a};



    \path[draw=black] (21.1374, 12.065) ellipse (0.8153cm and 0.8153cm);



    \node[anchor=south] (text129) at (21.1374, 11.9168){b.7};



    \path[draw=black] (18.8161, 12.1705).. controls (19.1477, 12.1553) and (19.5428, 12.137) .. (19.9108, 12.12);



    \path[draw=black,fill=black] (19.9094, 12.2435) -- (20.2561, 12.1042) -- (19.8981, 11.9969) -- (19.9094, 12.2435) -- (19.9094, 12.2435)-- cycle;;



    \node[anchor=south] (text3952) at (19.5636, 12.288){b};



    \path[draw=black] (21.9601, 12.065).. controls (22.2945, 12.065) and (22.691, 12.065) .. (23.0586, 12.065);



    \path[draw=black,fill=black] (23.0501, 12.1885) -- (23.4029, 12.065) -- (23.0501, 11.9415) -- (23.0501, 12.1885) -- (23.0501, 12.1885)-- cycle;;



    \node[anchor=south] (text4320) at (22.7111, 12.2132){a};



    \path[draw=black] (27.535, 13.7583) ellipse (0.8153cm and 0.8153cm);



    \path[draw=black] (27.535, 13.7583) ellipse (0.9564cm and 0.9564cm);



    \node[anchor=south] (text5) at (27.535, 13.6102){b.9};



    \path[draw=black] (24.9795, 12.4241).. controls (25.37, 12.6309) and (25.872, 12.8965) .. (26.325, 13.1364);



    \path[draw=black,fill=black] (26.258, 13.2408) -- (26.6273, 13.2965) -- (26.3733, 13.0224) -- (26.258, 13.2408) -- (26.258, 13.2408)-- cycle;;



    \node[anchor=south] (text3688) at (25.8202, 13.0598){b};



    \path[draw=black] (28.4974, 13.7252).. controls (29.7363, 13.6804) and (31.932, 13.6017) .. (33.2811, 13.5534);



    \path[draw=black,fill=black] (33.2691, 13.6772) -- (33.6173, 13.541) -- (33.2602, 13.4302) -- (33.2691, 13.6772) -- (33.2691, 13.6772)-- cycle;;



    \node[anchor=south] (text3593) at (31.0878, 13.8197){a};



    \path[draw=black] (26.5793, 13.9298).. controls (25.9422, 14.0335) and (25.073, 14.1464) .. (24.3011, 14.1464).. controls (17.974, 14.1464) and (17.974, 14.1464) .. (17.974, 14.1464).. controls (17.3704, 14.1464) and (16.6991, 14.1464) .. (16.1421, 14.1464);



    \path[draw=black,fill=black] (16.1477, 14.0229) -- (15.7949, 14.1464) -- (16.1477, 14.2699) -- (16.1477, 14.0229) -- (16.1477, 14.0229)-- cycle;;



    \node[anchor=south] (text8243) at (21.1374, 14.2946){b};



    \path[draw=black] (31.0878, 15.8397) ellipse (0.9663cm and 0.9663cm);



    \path[draw=black] (31.0878, 15.8397) ellipse (1.1074cm and 1.1074cm);



    \node[anchor=south] (text5084) at (31.0878, 15.6916){a.10};



    \path[draw=black] (28.3976, 14.2095).. controls (28.6981, 14.3768) and (29.0393, 14.5701) .. (29.3455, 14.7532).. controls (29.4954, 14.8428) and (29.6513, 14.9387) .. (29.8058, 15.0347);



    \path[draw=black,fill=black] (29.7275, 15.1317) -- (30.0919, 15.2153) -- (29.8595, 14.9225) -- (29.7275, 15.1317) -- (29.7275, 15.1317)-- cycle;;



    \node[anchor=south] (text9684) at (29.2358, 14.9013){a};



    \path[draw=black] (32.0135, 15.2157).. controls (32.4683, 14.8971) and (33.0221, 14.509) .. (33.486, 14.1838);



    \path[draw=black,fill=black] (33.5439, 14.2942) -- (33.7619, 13.9905) -- (33.4021, 14.0917) -- (33.5439, 14.2942) -- (33.5439, 14.2942)-- cycle;;



    \node[anchor=south] (text8658) at (32.9403, 14.7539){a};



    \path[draw=black] (29.9702, 15.8397).. controls (29.283, 15.8397) and (28.3743, 15.8397) .. (27.5703, 15.8397).. controls (17.974, 15.8397) and (17.974, 15.8397) .. (17.974, 15.8397).. controls (17.1782, 15.8397) and (16.3971, 15.3889) .. (15.827, 14.9525);



    \path[draw=black,fill=black] (15.9106, 14.8615) -- (15.5589, 14.7345) -- (15.7547, 15.053) -- (15.9106, 14.8615) -- (15.9106, 14.8615)-- cycle;;



    \node[anchor=south] (text7495) at (22.7111, 15.9879){b};



    \path[draw=black] (30.7411, 16.8945).. controls (30.7336, 17.2734) and (30.8494, 17.5821) .. (31.0878, 17.5821).. controls (31.2335, 17.5821) and (31.3334, 17.4674) .. (31.3873, 17.2921);



    \path[draw=black,fill=black] (31.5094, 17.3126) -- (31.4286, 16.9474) -- (31.2639, 17.2833) -- (31.5094, 17.3126) -- (31.5094, 17.3126)-- cycle;;



    \node[anchor=south] (text9014) at (31.0878, 17.7303){a};



  \end{scope}

\end{tikzpicture}

      \end{adjustbox}
      \caption{Автомат Глушкова}
      \label{fig:nfa:glushkov}
    \end{figure}

    Рассмотрим пары состояний $b.9$ и $a.10$, $S$ и $a.2$, $b.7$ и $a.4$. Для каждого состояния $q_i$ в паре запишем множество всех переходов, т.е. пар из состояния $q_j$ и символ перехода $\gamma$ таких, что $q_i \xrightarrow{\gamma} q_j$.

    %% Код простого reduce лежит в ветке `junk`.
    %% Псевдокод и теория для него взяты из учебника Javier Esparza "Automata theory. An algorithmic approach".

    \begin{enumerate}
      \itemsep -0.8pt
      \item $b.9$ и $a.10$:
        \begin{gather*}
          b.9 \to \{(b.5, b), (a.3, a), (a.10, a)\}\\
          a.10 \to \{(b.5, b), (a.3, a), (a.10, a)\}
        \end{gather*}
      \item $S$ и $a.2$:
        \begin{gather*}
          S \to \{(a.0, a), (b.5, b), (a.3, a)\}\\
          a.2 \to \{(a.0, a), (b.5, b), (a.3, a)\}
        \end{gather*}
      \item $b.7$ и $a.4$:
        \begin{gather*}
          b.7 \to \{(a.8, a)\}\\
          a.4 \to \{(a.8, a)\}
        \end{gather*}
    \end{enumerate}

    Заметим, что в каждой паре множество переходов полностью идентичны, поэтому состояния в паре лежат в одном классе эквивалентностей.

    Таким образом, получаем следующий НКА с 9 состояниями, также распознающий язык регулярного выражения $\re$.

    \begin{figure}[H]
      \centering
      \begin{adjustbox}{width=\textwidth}
      


\def \globalscale {1.000000}
\begin{tikzpicture}[y=1cm, x=1cm, yscale=\globalscale,xscale=\globalscale, every node/.append style={scale=\globalscale}, inner sep=0pt, outer sep=0pt]
  \begin{scope}[shift={(0.1411, -8.2247)}]
    \path[fill=white] (-0.1411, 8.2197) -- (-0.1411, 16.5855) -- (20.0741, 16.5855) -- (20.0741, 8.2197) -- (-0.1411, 8.2197) -- (-0.1411, 8.2197)-- cycle;;



    \path[draw=black] (2.2084, 11.3242) ellipse (0.635cm and 0.635cm);



    \path[draw=black] (2.2084, 11.3242) ellipse (0.7761cm and 0.7761cm);



    \node[anchor=south] (text9685) at (2.2084, 11.176){0};



    \path[draw=black] (5.1364, 9.9836) ellipse (0.635cm and 0.635cm);



    \node[anchor=south] (text247) at (5.1364, 9.8354){1};



    \path[draw=black] (2.921, 11.0063).. controls (3.296, 10.8303) and (3.768, 10.6091) .. (4.1765, 10.4172);



    \path[draw=black,fill=black] (4.2287, 10.5294) -- (4.4954, 10.2676) -- (4.124, 10.3057) -- (4.2287, 10.5294) -- (4.2287, 10.5294)-- cycle;;



    \node[anchor=south] (text3835) at (3.743, 10.8021){a};



    \path[draw=black] (5.1364, 13.7583) ellipse (0.635cm and 0.635cm);



    \node[anchor=south] (text2262) at (5.1364, 13.6102){3};



    \path[draw=black] (2.8219, 11.8166).. controls (3.254, 12.1846) and (3.8488, 12.6915) .. (4.3226, 13.0951);



    \path[draw=black,fill=black] (4.2386, 13.1858) -- (4.5872, 13.3205) -- (4.3988, 12.9977) -- (4.2386, 13.1858) -- (4.2386, 13.1858)-- cycle;;



    \node[anchor=south] (text249) at (3.743, 12.8351){b};



    \path[draw=black] (19.298, 12.8764) ellipse (0.635cm and 0.635cm);



    \node[anchor=south] (text9708) at (19.298, 12.7282){8};



    \path[draw=black] (2.5453, 10.6218).. controls (2.9743, 9.7455) and (3.8506, 8.3608) .. (5.1012, 8.3608).. controls (5.1012, 8.3608) and (5.1012, 8.3608) .. (16.4327, 8.3608).. controls (18.0721, 8.3608) and (18.7978, 10.5237) .. (19.0892, 11.8555);



    \path[draw=black,fill=black] (18.9653, 11.865) -- (19.1558, 12.1867) -- (19.2073, 11.8163) -- (18.9653, 11.865) -- (18.9653, 11.865)-- cycle;;



    \node[anchor=south] (text7550) at (10.7103, 8.509){a};



    \path[draw=black] (16.3975, 12.8764) ellipse (0.635cm and 0.635cm);



    \path[draw=black] (16.3975, 12.8764) ellipse (0.7761cm and 0.7761cm);



    \node[anchor=south] (text171) at (16.3975, 12.7282){7};



    \path[draw=black] (16.1195, 13.6169).. controls (16.0853, 13.9732) and (16.1777, 14.2875) .. (16.3975, 14.2875).. controls (16.528, 14.2875) and (16.6137, 14.1767) .. (16.6546, 14.012);



    \path[draw=black,fill=black] (16.7774, 14.0289) -- (16.673, 13.6701) -- (16.5308, 14.0159) -- (16.7774, 14.0289) -- (16.7774, 14.0289)-- cycle;;



    \node[anchor=south] (text1598) at (16.3975, 14.4357){a};



    \path[draw=black] (15.6104, 12.8333).. controls (15.0343, 12.8041) and (14.2198, 12.7706) .. (13.5047, 12.7706).. controls (7.8881, 12.7706) and (7.8881, 12.7706) .. (7.8881, 12.7706).. controls (7.2524, 12.7706) and (6.5804, 13.0126) .. (6.0639, 13.257);



    \path[draw=black,fill=black] (6.0089, 13.1466) -- (5.7503, 13.4161) -- (6.1207, 13.3667) -- (6.0089, 13.1466) -- (6.0089, 13.1466)-- cycle;;



    \node[anchor=south] (text150) at (10.7103, 12.9187){b};



    \path[draw=black] (17.1912, 12.8764).. controls (17.5147, 12.8764) and (17.8971, 12.8764) .. (18.245, 12.8764);



    \path[draw=black,fill=black] (18.239, 12.9999) -- (18.5917, 12.8764) -- (18.239, 12.7529) -- (18.239, 12.9999) -- (18.239, 12.9999)-- cycle;;



    \node[anchor=south] (text733) at (17.9183, 13.0246){a};



    \path[draw=black,fill=black] (0.0635, 11.3242) ellipse (0.0635cm and 0.0635cm);



    \path[draw=black] (0.1348, 11.3242).. controls (0.2501, 11.3242) and (0.6265, 11.3242) .. (1.0252, 11.3242);



    \path[draw=black,fill=black] (1.016, 11.4476) -- (1.3688, 11.3242) -- (1.016, 11.2007) -- (1.016, 11.4476) -- (1.016, 11.4476)-- cycle;;



    \path[draw=black] (7.9234, 10.5833) ellipse (0.635cm and 0.635cm);



    \node[anchor=south] (text4268) at (7.9234, 10.4352){2};



    \path[draw=black] (5.7746, 10.1166).. controls (6.1052, 10.1896) and (6.5232, 10.2821) .. (6.9003, 10.3653);



    \path[draw=black,fill=black] (6.8636, 10.4835) -- (7.2348, 10.439) -- (6.9169, 10.2426) -- (6.8636, 10.4835) -- (6.8636, 10.4835)-- cycle;;



    \node[anchor=south] (text663) at (6.5299, 10.4496){b};



    \path[draw=black] (7.9234, 14.4992) ellipse (0.635cm and 0.635cm);



    \node[anchor=south] (text3962) at (7.9234, 14.351){4};



    \path[draw=black] (5.7612, 13.9192).. controls (6.0992, 14.0113) and (6.5324, 14.1295) .. (6.919, 14.2349);



    \path[draw=black,fill=black] (6.8756, 14.351) -- (7.2482, 14.3245) -- (6.9405, 14.1129) -- (6.8756, 14.351) -- (6.8756, 14.351)-- cycle;;



    \node[anchor=south] (text9652) at (6.5299, 14.2991){a};



    \path[draw=black] (10.7103, 14.6403) ellipse (0.635cm and 0.635cm);



    \node[anchor=south] (text4369) at (10.7103, 14.4921){5};



    \path[draw=black] (18.9526, 13.425).. controls (18.6016, 13.9852) and (17.9737, 14.8255) .. (17.1736, 15.2047).. controls (15.088, 16.1936) and (14.1997, 15.9103) .. (11.9803, 15.2753).. controls (11.8614, 15.2411) and (11.7404, 15.1952) .. (11.623, 15.1437);



    \path[draw=black,fill=black] (11.6861, 15.0372) -- (11.3153, 14.9902) -- (11.5761, 15.258) -- (11.6861, 15.0372) -- (11.6861, 15.0372)-- cycle;;



    \node[anchor=south] (text7203) at (14.8629, 15.9999){a};



    \path[draw=black] (7.3187, 10.826).. controls (7.1124, 10.9054) and (6.8756, 10.9873) .. (6.6534, 11.0419).. controls (5.3351, 11.3658) and (4.9766, 11.4296) .. (3.6195, 11.3947).. controls (3.5472, 11.393) and (3.4731, 11.3905) .. (3.3983, 11.3877);



    \path[draw=black,fill=black] (3.4117, 11.2645) -- (3.0536, 11.3725) -- (3.4004, 11.5115) -- (3.4117, 11.2645) -- (3.4117, 11.2645)-- cycle;;



    \node[anchor=south] (text8989) at (5.1364, 11.5468){a};



    \path[draw=black] (8.575, 14.5313).. controls (8.8946, 14.5479) and (9.2932, 14.5687) .. (9.6569, 14.5874);



    \path[draw=black,fill=black] (9.6492, 14.7105) -- (10.008, 14.6057) -- (9.6619, 14.4639) -- (9.6492, 14.7105) -- (9.6492, 14.7105)-- cycle;;



    \node[anchor=south] (text6999) at (9.3169, 14.7221){b};



    \path[draw=black] (13.4694, 14.3933) ellipse (0.635cm and 0.635cm);



    \node[anchor=south] (text6552) at (13.4694, 14.2452){6};



    \path[draw=black] (11.3556, 14.5842).. controls (11.672, 14.5553) and (12.0664, 14.5189) .. (12.4266, 14.4858);



    \path[draw=black,fill=black] (12.4333, 14.6092) -- (12.7734, 14.454) -- (12.4107, 14.3633) -- (12.4333, 14.6092) -- (12.4333, 14.6092)-- cycle;;



    \node[anchor=south] (text8189) at (12.09, 14.6724){a};



    \path[draw=black] (14.0423, 14.1079).. controls (14.4085, 13.9136) and (14.9049, 13.65) .. (15.3448, 13.4165);



    \path[draw=black,fill=black] (15.3956, 13.5294) -- (15.6492, 13.2549) -- (15.2799, 13.3114) -- (15.3956, 13.5294) -- (15.3956, 13.5294)-- cycle;;



    \node[anchor=south] (text8939) at (14.8629, 13.8737){b};



  \end{scope}

\end{tikzpicture}

      \end{adjustbox}
      \caption{Редуцированный автомат Глушкова}
      \label{fig:nfa:glushkov-reduced}
    \end{figure}

    \begin{figure}[H]
      \centering
      \begin{tabular}{|c|c|c|c|c|c|c|}
        \hline
        & aba & abab & aab & a & bab & b \\
        \hline
        $\varepsilon$ & $+$ & $-$ & $-$ & $-$ & $-$ & $-$ \\
        \hline
        b & $-$ & $+$ & $-$ & $-$ & $-$ & $-$ \\
        \hline
        a & $-$ & $-$ & $+$ & $-$ & $-$ & $-$ \\
        \hline
        ab & $-$ & $-$ & $-$ & $+$ & $-$ & $-$ \\
        \hline
        ba & $-$ & $-$ & $-$ & $-$ & $+$ & $-$ \\
        \hline
        aaa & $-$ & $-$ & $-$ & $-$ & $-$ & $+$\\
        \hline
      \end{tabular}
      \caption{Таблица префиксов-суффиксов}
      \label{table:nfa-prefix-suffix}
    \end{figure}

    \section{Детерминированный конечный автомат}

    %% Строим через Chipollino.
    %% Minimize.Glushkov {(aba)*((aa|bab)aba*)*} !!

    %% Можно было бы сделать автомат и с ловушкой, но он бы выглядел сильно хуже
    %% из-за большого числа дополнительных связей.

    \begin{figure}[H]
      \centering
      \begin{adjustbox}{width=\textwidth}
      

\def \globalscale {1.000000}
\begin{tikzpicture}[y=1cm, x=1cm, yscale=\globalscale,xscale=\globalscale, every node/.append style={scale=\globalscale}, inner sep=0pt, outer sep=0pt]
  \begin{scope}[shift={(0.1411, -14.0226)}]
    \path[fill=white] (-0.1411, 14.0053) -- (-0.1411, 28.169) -- (45.5923, 28.169) -- (45.5923, 14.0053) -- (-0.1411, 14.0053) -- (-0.1411, 14.0053)-- cycle;;



    \path[draw=black] (2.2084, 25.8477) ellipse (0.635cm and 0.635cm);



    \path[draw=black] (2.2084, 25.8477) ellipse (0.7761cm and 0.7761cm);



    \node[anchor=south] (text6664) at (2.2084, 25.6995){0};



    \path[draw=black] (5.1086, 23.9427) ellipse (0.635cm and 0.635cm);



    \node[anchor=south] (text9621) at (5.1086, 23.7945){1};



    \path[draw=black] (2.8176, 25.3619).. controls (3.0589, 25.1661) and (3.3464, 24.9442) .. (3.6195, 24.7611).. controls (3.7955, 24.6433) and (3.9896, 24.5258) .. (4.1769, 24.4186);



    \path[draw=black,fill=black] (4.2277, 24.5315) -- (4.476, 24.2524) -- (4.1077, 24.3156) -- (4.2277, 24.5315) -- (4.2277, 24.5315)-- cycle;;



    \node[anchor=south] (text1977) at (3.7292, 24.9093){a};



    \path[draw=black] (41.8207, 22.3905) ellipse (0.635cm and 0.635cm);



    \node[anchor=south] (text8861) at (41.8207, 22.2423){3};



    \path[draw=black] (2.7891, 26.3765).. controls (3.3249, 26.8383) and (4.1889, 27.4352) .. (5.0733, 27.4352).. controls (5.0733, 27.4352) and (5.0733, 27.4352) .. (38.5967, 27.4352).. controls (40.4774, 27.4352) and (41.2944, 24.8994) .. (41.6105, 23.4234);



    \path[draw=black,fill=black] (41.7304, 23.452) -- (41.6779, 23.0819) -- (41.4884, 23.404) -- (41.7304, 23.452) -- (41.7304, 23.452)-- cycle;;



    \node[anchor=south] (text4115) at (21.2005, 27.5833){b};



    \path[draw=black] (13.583, 23.5899) ellipse (0.635cm and 0.635cm);



    \path[draw=black] (13.583, 23.5899) ellipse (0.7761cm and 0.7761cm);



    \node[anchor=south] (text9948) at (13.583, 23.4417){7};



    \path[draw=black] (16.6243, 22.2141) ellipse (0.635cm and 0.635cm);



    \path[draw=black] (16.6243, 22.2141) ellipse (0.7761cm and 0.7761cm);



    \node[anchor=south] (text3039) at (16.6243, 22.0659){8};



    \path[draw=black] (14.2854, 23.2396).. controls (14.509, 23.1253) and (14.7609, 23.0001) .. (14.9941, 22.8914).. controls (15.1613, 22.8138) and (15.3398, 22.7344) .. (15.5145, 22.6586);



    \path[draw=black,fill=black] (15.5603, 22.7732) -- (15.8365, 22.5217) -- (15.4637, 22.546) -- (15.5603, 22.7732) -- (15.5603, 22.7732)-- cycle;;



    \node[anchor=south] (text8985) at (15.1035, 23.0396){a};



    \path[draw=black] (14.3595, 23.7412).. controls (14.9585, 23.8502) and (15.8246, 23.978) .. (16.589, 23.978).. controls (16.589, 23.978) and (16.589, 23.978) .. (38.5967, 23.978).. controls (39.4892, 23.978) and (40.3885, 23.4738) .. (41.0034, 23.0336);



    \path[draw=black,fill=black] (41.064, 23.1426) -- (41.2715, 22.8314) -- (40.9155, 22.9454) -- (41.064, 23.1426) -- (41.064, 23.1426)-- cycle;;



    \node[anchor=south] (text1122) at (27.7188, 24.1261){b};



    \path[draw=black] (19.666, 19.4271) ellipse (0.635cm and 0.635cm);



    \path[draw=black] (19.666, 19.4271) ellipse (0.7761cm and 0.7761cm);



    \node[anchor=south] (text2916) at (19.666, 19.279){9};



    \path[draw=black] (17.2184, 21.6898).. controls (17.6583, 21.2774) and (18.2785, 20.6954) .. (18.7801, 20.2251);



    \path[draw=black,fill=black] (18.8623, 20.3172) -- (19.0352, 19.9856) -- (18.6933, 20.1369) -- (18.8623, 20.3172) -- (18.8623, 20.3172)-- cycle;;



    \node[anchor=south] (text5065) at (18.1451, 21.0524){a};



    \path[draw=black] (17.4036, 22.2829).. controls (18.004, 22.3322) and (18.8711, 22.3905) .. (19.6307, 22.3905).. controls (19.6307, 22.3905) and (19.6307, 22.3905) .. (38.5967, 22.3905).. controls (39.3255, 22.3905) and (40.1479, 22.3905) .. (40.7748, 22.3905);



    \path[draw=black,fill=black] (40.7681, 22.5139) -- (41.1208, 22.3905) -- (40.7681, 22.267) -- (40.7681, 22.5139) -- (40.7681, 22.5139)-- cycle;;



    \node[anchor=south] (text4411) at (29.3483, 22.5386){b};



    \path[draw=black] (22.83, 15.0174) ellipse (0.7299cm and 0.7299cm);



    \path[draw=black] (22.83, 15.0174) ellipse (0.871cm and 0.871cm);



    \node[anchor=south] (text2458) at (22.83, 14.8692){10};



    \path[draw=black] (20.0459, 18.7357).. controls (20.3055, 18.2376) and (20.683, 17.5581) .. (21.0771, 17.0).. controls (21.3356, 16.6335) and (21.6496, 16.2556) .. (21.9382, 15.9293);



    \path[draw=black,fill=black] (22.0239, 16.0186) -- (22.1682, 15.6739) -- (21.8405, 15.8531) -- (22.0239, 16.0186) -- (22.0239, 16.0186)-- cycle;;



    \node[anchor=south] (text5606) at (21.2005, 17.1482){a};



    \path[draw=black] (20.3189, 19.8737).. controls (20.9204, 20.2636) and (21.8758, 20.7677) .. (22.7947, 20.7677).. controls (22.7947, 20.7677) and (22.7947, 20.7677) .. (27.7541, 20.7677).. controls (29.833, 20.7677) and (30.3523, 20.803) .. (32.4312, 20.803).. controls (32.4312, 20.803) and (32.4312, 20.803) .. (38.5967, 20.803).. controls (39.4892, 20.803) and (40.3885, 21.3074) .. (41.0034, 21.7473);



    \path[draw=black,fill=black] (40.9155, 21.8359) -- (41.2715, 21.9495) -- (41.064, 21.6387) -- (40.9155, 21.8359) -- (40.9155, 21.8359)-- cycle;;



    \node[anchor=south] (text4994) at (30.9781, 20.9455){b};



    \path[draw=black] (22.5291, 15.8387).. controls (22.5005, 16.2059) and (22.6007, 16.5238) .. (22.83, 16.5238).. controls (22.9697, 16.5238) and (23.0614, 16.4056) .. (23.1052, 16.2302);



    \path[draw=black,fill=black] (23.2276, 16.2518) -- (23.1274, 15.8919) -- (22.9814, 16.2359) -- (23.2276, 16.2518) -- (23.2276, 16.2518)-- cycle;;



    \node[anchor=south] (text9190) at (22.83, 16.6719){a};



    \path[draw=black] (26.0893, 15.6877) ellipse (0.7299cm and 0.7299cm);



    \path[draw=black] (26.0893, 15.6877) ellipse (0.871cm and 0.871cm);



    \node[anchor=south] (text3094) at (26.0893, 15.5395){11};



    \path[draw=black] (23.6866, 15.1903).. controls (24.0369, 15.264) and (24.4535, 15.3515) .. (24.8384, 15.4323);



    \path[draw=black,fill=black] (24.8017, 15.5508) -- (25.1721, 15.5025) -- (24.8525, 15.3091) -- (24.8017, 15.5508) -- (24.8017, 15.5508)-- cycle;;



    \node[anchor=south] (text5953) at (24.4595, 15.5208){b};



    \path[draw=black] (29.3483, 16.1463) ellipse (0.7299cm and 0.7299cm);



    \path[draw=black] (29.3483, 16.1463) ellipse (0.871cm and 0.871cm);



    \node[anchor=south] (text662) at (29.3483, 15.9981){12};



    \path[draw=black] (26.9621, 15.8083).. controls (27.3036, 15.8574) and (27.7054, 15.9152) .. (28.0793, 15.9688);



    \path[draw=black,fill=black] (28.0501, 16.0895) -- (28.417, 16.0175) -- (28.0853, 15.845) -- (28.0501, 16.0895) -- (28.0501, 16.0895)-- cycle;;



    \node[anchor=south] (text1609) at (27.7188, 16.0789){a};



    \path[draw=black] (26.7755, 16.2405).. controls (27.2274, 16.6) and (27.8578, 17.0543) .. (28.4773, 17.3457).. controls (29.5049, 17.829) and (29.8073, 18.016) .. (30.9428, 18.016).. controls (30.9428, 18.016) and (30.9428, 18.016) .. (38.5967, 18.016).. controls (40.2505, 18.016) and (41.1374, 20.0826) .. (41.5301, 21.3752);



    \path[draw=black,fill=black] (41.4108, 21.4069) -- (41.626, 21.7124) -- (41.6486, 21.3395) -- (41.4108, 21.4069) -- (41.4108, 21.4069)-- cycle;;



    \node[anchor=south] (text4263) at (33.9549, 18.1642){b};



    \path[draw=black] (28.5203, 16.4356).. controls (28.3012, 16.5142) and (28.0628, 16.5993) .. (27.8423, 16.6755).. controls (24.8472, 17.714) and (24.0721, 17.8961) .. (21.0771, 18.9332).. controls (20.9825, 18.966) and (20.8844, 19.0003) .. (20.7864, 19.0352);



    \path[draw=black,fill=black] (20.7553, 18.9152) -- (20.4643, 19.1498) -- (20.8382, 19.1477) -- (20.7553, 18.9152) -- (20.7553, 18.9152)-- cycle;;



    \node[anchor=south] (text2553) at (24.4595, 17.9811){a};



    \path[draw=black] (32.4665, 16.1463) ellipse (0.7299cm and 0.7299cm);



    \node[anchor=south] (text9625) at (32.4665, 15.9981){13};



    \path[draw=black] (30.2327, 16.1463).. controls (30.5693, 16.1463) and (30.9612, 16.1463) .. (31.32, 16.1463);



    \path[draw=black,fill=black] (31.3147, 16.2698) -- (31.6674, 16.1463) -- (31.3147, 16.0228) -- (31.3147, 16.2698) -- (31.3147, 16.2698)-- cycle;;



    \node[anchor=south] (text6105) at (30.9781, 16.2945){b};



    \path[draw=black] (38.5614, 16.1463) ellipse (0.7299cm and 0.7299cm);



    \path[draw=black] (38.5614, 16.1463) ellipse (0.871cm and 0.871cm);



    \node[anchor=south] (text9734) at (38.5614, 15.9981){15};



    \path[draw=black] (41.8207, 15.476) ellipse (0.7299cm and 0.7299cm);



    \path[draw=black] (41.8207, 15.476) ellipse (0.871cm and 0.871cm);



    \node[anchor=south] (text1479) at (41.8207, 15.3278){16};



    \path[draw=black] (39.3573, 15.7572).. controls (39.5806, 15.6609) and (39.829, 15.5702) .. (40.0678, 15.5183).. controls (40.2209, 15.4852) and (40.3832, 15.4637) .. (40.5448, 15.4506);



    \path[draw=black,fill=black] (40.5388, 15.5748) -- (40.8859, 15.4361) -- (40.5282, 15.3278) -- (40.5388, 15.5748) -- (40.5388, 15.5748)-- cycle;;



    \node[anchor=south] (text3483) at (40.1913, 15.6665){a};



    \path[draw=black] (39.4169, 16.3548).. controls (39.7379, 16.4747) and (40.0837, 16.6571) .. (40.3147, 16.9295).. controls (41.3833, 18.1903) and (41.6825, 20.1542) .. (41.7622, 21.3515);



    \path[draw=black,fill=black] (41.6387, 21.3558) -- (41.7798, 21.7018) -- (41.8853, 21.3431) -- (41.6387, 21.3558) -- (41.6387, 21.3558)-- cycle;;



    \node[anchor=south] (text4910) at (40.1913, 17.0776){b};



    \path[draw=black] (41.5265, 16.3064).. controls (41.1103, 17.4562) and (40.175, 19.3918) .. (38.5967, 19.3918).. controls (32.4312, 19.3918) and (32.4312, 19.3918) .. (32.4312, 19.3918).. controls (30.3523, 19.3918) and (29.833, 19.4271) .. (27.7541, 19.4271).. controls (22.7947, 19.4271) and (22.7947, 19.4271) .. (22.7947, 19.4271).. controls (22.1506, 19.4271) and (21.4313, 19.4271) .. (20.8485, 19.4271);



    \path[draw=black,fill=black] (20.8573, 19.3036) -- (20.5045, 19.4271) -- (20.8573, 19.5506) -- (20.8573, 19.3036) -- (20.8573, 19.3036)-- cycle;;



    \node[anchor=south] (text3518) at (30.9781, 19.5478){a};



    \path[draw=black] (41.0545, 15.0294).. controls (40.4262, 14.6886) and (39.4839, 14.2766) .. (38.5967, 14.2766).. controls (29.313, 14.2766) and (29.313, 14.2766) .. (29.313, 14.2766).. controls (28.539, 14.2766) and (27.7424, 14.629) .. (27.1392, 14.9814);



    \path[draw=black,fill=black] (27.0884, 14.8675) -- (26.8531, 15.1582) -- (27.2182, 15.0777) -- (27.0884, 14.8675) -- (27.0884, 14.8675)-- cycle;;



    \node[anchor=south] (text1686) at (33.9549, 14.4247){b};



    \path[draw=black,fill=black] (0.0635, 25.8477) ellipse (0.0635cm and 0.0635cm);



    \path[draw=black] (0.1348, 25.8477).. controls (0.2501, 25.8477) and (0.6265, 25.8477) .. (1.0252, 25.8477);



    \path[draw=black,fill=black] (1.016, 25.9711) -- (1.3688, 25.8477) -- (1.016, 25.7242) -- (1.016, 25.9711) -- (1.016, 25.9711)-- cycle;;



    \path[draw=black] (7.8955, 25.6713) ellipse (0.635cm and 0.635cm);



    \node[anchor=south] (text4078) at (7.8955, 25.5231){2};



    \path[draw=black] (5.6889, 24.2217).. controls (5.9732, 24.3727) and (6.3239, 24.5678) .. (6.6255, 24.7611).. controls (6.7695, 24.8535) and (6.9187, 24.9566) .. (7.0616, 25.0592);



    \path[draw=black,fill=black] (6.9779, 25.1509) -- (7.335, 25.2614) -- (7.1247, 24.9523) -- (6.9779, 25.1509) -- (6.9779, 25.1509)-- cycle;;



    \node[anchor=south] (text6733) at (6.502, 24.9093){b};



    \path[draw=black] (7.8955, 23.5899) ellipse (0.635cm and 0.635cm);



    \node[anchor=south] (text341) at (7.8955, 23.4417){5};



    \path[draw=black] (5.666, 23.6283).. controls (5.879, 23.5172) and (6.133, 23.4061) .. (6.3786, 23.35).. controls (6.5398, 23.3133) and (6.7123, 23.3098) .. (6.8802, 23.3246);



    \path[draw=black,fill=black] (6.8591, 23.4463) -- (7.2274, 23.3842) -- (6.9007, 23.2029) -- (6.8591, 23.4463) -- (6.8591, 23.4463)-- cycle;;



    \node[anchor=south] (text8895) at (6.502, 23.4982){a};



    \path[draw=black] (44.8162, 24.0132) ellipse (0.635cm and 0.635cm);



    \node[anchor=south] (text2563) at (44.8162, 23.8651){4};



    \path[draw=black] (42.3926, 22.6882).. controls (42.8103, 22.92) and (43.3998, 23.247) .. (43.8884, 23.5183);



    \path[draw=black,fill=black] (43.8228, 23.6231) -- (44.1911, 23.6862) -- (43.9427, 23.4072) -- (43.8228, 23.6231) -- (43.8228, 23.6231)-- cycle;;



    \node[anchor=south] (text7091) at (43.4365, 23.4474){a};



    \path[draw=black] (7.2524, 25.6903).. controls (6.3327, 25.7193) and (4.572, 25.7747) .. (3.4025, 25.8113);



    \path[draw=black,fill=black] (3.4032, 25.6879) -- (3.0547, 25.8223) -- (3.411, 25.9348) -- (3.4032, 25.6879) -- (3.4032, 25.6879)-- cycle;;



    \node[anchor=south] (text2271) at (5.1086, 25.9253){a};



    \path[draw=black] (10.6549, 23.5899) ellipse (0.635cm and 0.635cm);



    \node[anchor=south] (text6704) at (10.6549, 23.4417){6};



    \path[draw=black] (8.5411, 23.5899).. controls (8.8551, 23.5899) and (9.2463, 23.5899) .. (9.604, 23.5899);



    \path[draw=black,fill=black] (9.5963, 23.7134) -- (9.949, 23.5899) -- (9.5963, 23.4664) -- (9.5963, 23.7134) -- (9.5963, 23.7134)-- cycle;;



    \node[anchor=south] (text8572) at (9.2752, 23.7381){a};



    \path[draw=black] (44.3353, 24.4387).. controls (43.7942, 24.9089) and (42.8325, 25.6007) .. (41.856, 25.6007).. controls (10.6197, 25.6007) and (10.6197, 25.6007) .. (10.6197, 25.6007).. controls (9.7737, 25.6007) and (9.0226, 24.9548) .. (8.5312, 24.3949);



    \path[draw=black,fill=black] (8.6392, 24.3321) -- (8.3196, 24.1385) -- (8.449, 24.4891) -- (8.6392, 24.3321) -- (8.6392, 24.3321)-- cycle;;



    \node[anchor=south] (text923) at (26.0893, 25.7489){b};



    \path[draw=black] (11.2963, 23.5899).. controls (11.6127, 23.5899) and (12.011, 23.5899) .. (12.3846, 23.5899);



    \path[draw=black,fill=black] (12.3836, 23.7134) -- (12.7363, 23.5899) -- (12.3836, 23.4664) -- (12.3836, 23.7134) -- (12.3836, 23.7134)-- cycle;;



    \node[anchor=south] (text5196) at (12.0484, 23.7381){b};



    \path[draw=black] (35.4436, 16.1463) ellipse (0.7299cm and 0.7299cm);



    \node[anchor=south] (text3376) at (35.4436, 15.9981){14};



    \path[draw=black] (33.2052, 16.1463).. controls (33.5351, 16.1463) and (33.9355, 16.1463) .. (34.3038, 16.1463);



    \path[draw=black,fill=black] (34.2949, 16.2698) -- (34.6477, 16.1463) -- (34.2949, 16.0228) -- (34.2949, 16.2698) -- (34.2949, 16.2698)-- cycle;;



    \node[anchor=south] (text1679) at (33.9549, 16.2945){a};



    \path[draw=black] (36.1855, 16.1463).. controls (36.51, 16.1463) and (36.9048, 16.1463) .. (37.277, 16.1463);



    \path[draw=black,fill=black] (37.2763, 16.2698) -- (37.629, 16.1463) -- (37.2763, 16.0228) -- (37.2763, 16.2698) -- (37.2763, 16.2698)-- cycle;;



    \node[anchor=south] (text9859) at (36.932, 16.2945){b};



  \end{scope}

\end{tikzpicture}

      \end{adjustbox}
      \caption{Минимальный ДКА}
      \label{fig:dfa}
    \end{figure}

    Проверим минимальность. Для каждой пары состояний запишем минимальное $k$ такое, что $q_i \not \sim_{k} q_j$.

    %% Таблицу тоже генерировал кодом, вручную сложно сравнивать такое количество состояний,
    %% а просто сказать про различимость (или поставить везде галочки) -- как-то неправильно.

    \begin{figure}[H]
      \centering
      \begin{tabular}{|>{\columncolor{gray!25}}c|c|c|c|c|c|c|c|c|c|c|c|c|c|c|c|c|c|c|}
      \hline
      \rowcolor{gray!45} & 0 & 1 & 2 & 3 & 4 & 5 & 6 & 7 & 8 & 9 & 10 & 11 & 12 & 13 & 14 & 15 & 16 & T \\
      \hline
      0  & = &   &   &   &   &   &   &   &   &   &   &   &   &   &   &   &   &  \\
      \hline
      1  & 0 & = &   &   &   &   &   &   &   &   &   &   &   &   &   &   &   &  \\
      \hline
      2  & 0 & 1 & = &   &   &   &   &   &   &   &   &   &   &   &   &   &   &  \\
      \hline
      3  & 0 & 3 & 1 & = &   &   &   &   &   &   &   &   &   &   &   &   &   &  \\
      \hline
      4  & 0 & 3 & 1 & 4 & = &   &   &   &   &   &   &   &   &   &   &   &   &  \\
      \hline
      5  & 0 & 3 & 1 & 2 & 2 & = &   &   &   &   &   &   &   &   &   &   &   &  \\
      \hline
      6  & 0 & 1 & 1 & 1 & 1 & 2 & = &   &   &   &   &   &   &   &   &   &   &  \\
      \hline
      7  & 1 & 0 & 0 & 0 & 0 & 0 & 0 & = &   &   &   &   &   &   &   &   &   &  \\
      \hline
      8  & 1 & 0 & 0 & 0 & 0 & 0 & 0 & 3 & = &   &   &   &   &   &   &   &   &  \\
      \hline
      9  & 1 & 0 & 0 & 0 & 0 & 0 & 0 & 2 & 2 & = &   &   &   &   &   &   &   &  \\
      \hline
      10 & 1 & 0 & 0 & 0 & 0 & 0 & 0 & 3 & 2 & 1 & = &   &   &   &   &   &   &  \\
      \hline
      11 & 1 & 0 & 0 & 0 & 0 & 0 & 0 & 4 & 3 & 2 & 3 & = &   &   &   &   &   &  \\
      \hline
      12 & 1 & 0 & 0 & 0 & 0 & 0 & 0 & 3 & 3 & 2 & 2 & 3 & = &   &   &   &   &  \\
      \hline
      13 & 0 & 3 & 1 & 2 & 2 & 4 & 1 & 0 & 0 & 0 & 0 & 0 & 0 & = &   &   &   &  \\
      \hline
      14 & 0 & 3 & 1 & 4 & 1 & 2 & 3 & 0 & 0 & 0 & 0 & 0 & 0 & 2 & = &   &   &  \\
      \hline
      15 & 1 & 0 & 0 & 0 & 0 & 0 & 0 & 2 & 2 & 3 & 3 & 2 & 2 & 0 & 0 & = &   &  \\
      \hline
      16 & 1 & 0 & 0 & 0 & 0 & 0 & 0 & 1 & 1 & 1 & 2 & 1 & 1 & 0 & 0 & 2 & = &  \\
      \hline
      T  & 0 & 3 & 1 & 4 & 3 & 2 & 1 & 0 & 0 & 0 & 0 & 0 & 0 & 2 & 1 & 0 & 0 & =\\
      \hline
      \end{tabular}
      \caption{Таблица различимости состояний}
      \label{table:equivalence-classes}
    \end{figure}

    Из таблицы \ref{table:equivalence-classes} видно, что все состояния попарно различимы, следовательно, ДКА, представленный на рисунке \ref{fig:dfa}, является минимальным.

    \section{Переключающийся конечный автомат}

    %% Шизофазия будущего академщика:
    %% -----------------------------
    %% tl;dr что-то нормальное не осилил
    %%
    %% ПКА --- самая грустная тема в этой лабе, так как здесь буквально негде использовать
    %% ПКА наполную. Значительную часть времени сложнее понять не сам ПКА, а можно ли построить
    %% его лучше, чем простым переписыванием НКА, т.е., чтобы у него было меньше состояний, чем
    %% в НКА (в автомате Глушкова или Томпсона).
    %%
    %% Понятно, конечно, что можно сгенерировать все возможные ПКА и проверить, если они
    %% соответствуют исходной регуляке, но это уже извращение, кмк, да и не является каким-то 
    %% интересным (или, честно говоря, полезным) решением.
    %%
    %% Логически, основной плюс использования ПКА (для минимизации количества состояний),
    %% возникает, когда задачу можно разбить на несколько небольших подзадач, которые можно
    %% проверять "параллельно" и независимо друг от друга. Самый заезженный пример (был даже на
    %% одном из семинаров) --- делимость на 6. Тут, очевидно, что 6 = 3 x 2, а значит можно
    %% проверять только делимость на 2 и на 3, причём проверки не влияют друг на друга, поэтому 
    %% их можно производить  "параллельно" (т.е., в отдельных НКА) и потом объединять результат
    %% вычислений через логическое И. Но такой трюк, очевидно, нельзя применить для случая
    %% делимости на 4, потому что там проверки зависят друг от друга, а значит вычисляются не
    %% "параллельно".
    %%
    %% Аналогично можно посмотреть на регулярку --- здесь состояния внутри одного блока зависят
    %% друг от друга... Поэтому параллельность должна строиться над блоками или совокупностью
    %% блоков. Но тут появляется логическая проблема с тем, что непонятно как найти такие $r_1$
    %% и $r_2$, что пересечение регулярок $r_1$ и $r_2$ совпадало с самой регуляркой $r$, причём
    %% $|Nfa_min(r_1)| + |Nfa_min(r_2)| < |Nfa_min(r)|$ и возможно ли это? Ощущение, что есть,
    %% наверняка, какие-то частные (или, возможно, наоборот, частые) случаи, когда такие $r_1$
    %% и $r_2$ находятся с прочтения, но в голову ничего подобного не приходит. Поэтому "лучшее",
    %% что могу предложить (для исходной регулярки) --- ПКА, который пересекает НКА по
    %% префиксу/суффиксу и НКА на весь язык. Понятное дело, что это бессмыслная работа, но,
    %% к сожалению, ничего умнее в голову не приходит, сколько её не ломай. Делить на какие-то
    %% блоки не получится, потому что, опять таки, теряется независимость между вычисленями и не
    %% получается "распареллелить", а значит и "сэкономить" состояний.
    %%
    %% Понятное дело, что это всё имхо, собственные догадки и попытки осознать, и в реальности
    %% наверняка есть какая-то теорема/алгоритм/научная статья, которые не увидел, и которые
    %% полностью объясняют все непонятные моменты и предлагают элегантные методы построения
    %% подобных структур... Но что поделать, работаю с тем, что есть.

    %% Хитрость, которую не реализовал (небольшая выжимка из [1]):
    %% ----------------------------------------------------------
    %%
    %% Введём особый тип ПКА. s-ПКА это ПКА, где начальное состояние не является частью
    %% какого-либо цикла в автомате. Очевидно, что для каждого ПКА с k состояниями существует
    %% s-ПКА, количество состояний которого равны k+1. Обратное не всегда справедливо.
    %%
    %% Согласно [1, Corollary 7.6] если в минимальном ДКА, который распознаёт язык L^{R}, всего
    %% n состояний, то у соответствующего ему s-ПКА будет 1 + ceil(log2(n)) состояний. А значит
    %% минимальный ПКА может иметь либо 1 + ceil(log2(n)) состояний, либо ceil(log2(n)) состояний.
    %%
    %% Для исходного регулярного выражения размер минимального ДКА, который распознаёт L^{R},
    %% равен 16 (считал через Chipollino) с учётом состояния ловушки. А значит размер минимального
    %% ПКА равен либо 4, либо 5 состояниям.
    %%
    %% [1] A. Fellah, H. Jürgensen, and S. Yu. 1990. Constructions for alternating finite automata *.
    %%     International Journal of Computer Mathematics 35, 1–4 (1990), 117–132.
    %%     DOI:https://doi.org/10.1080/00207169008803893

    \begin{figure}[H]
      \centering
      \begin{adjustbox}{width=\textwidth}
      \def \globalscale {1.000000}
\begin{tikzpicture}[y=1cm, x=1cm, yscale=\globalscale,xscale=\globalscale, every node/.append style={scale=\globalscale}, inner sep=0pt, outer sep=0pt]
  \begin{scope}[shift={(0.1411, -14.1164)}]
    \path[fill=white] (-0.1411, 14.1111) -- (-0.1411, 28.3686) -- (26.5875, 28.3686) -- (26.5875, 14.1111) -- (-0.1411, 14.1111) -- (-0.1411, 14.1111)-- cycle;;



    \path[draw=black] (2.2278, 20.2494) ellipse (0.6544cm and 0.6544cm);



    \path[draw=black] (2.2278, 20.2494) ellipse (0.7955cm and 0.7955cm);



    \node[anchor=south] (text324) at (2.2278, 20.1013){\&};



    \path[draw=black] (6.7423, 23.5303) ellipse (0.8812cm and 0.8812cm);



    \path[draw=black] (6.7423, 23.5303) ellipse (1.0224cm and 1.0224cm);



    \node[anchor=south] (text359) at (6.7423, 23.3821){$a_1$};



    \path[draw=black] (2.8903, 20.7126).. controls (3.5874, 21.227) and (4.7269, 22.0684) .. (5.5795, 22.6981);



    \path[draw=black,fill=black] (5.5023, 22.7944) -- (5.8593, 22.9045) -- (5.6487, 22.5958) -- (5.5023, 22.7944) -- (5.5023, 22.7944)-- cycle;;



    \node[anchor=south] (text4387) at (4.3716, 22.4282){$\varepsilon$};



    \path[draw=black] (6.7423, 18.3797) ellipse (0.635cm and 0.635cm);



    \path[draw=black] (6.7423, 18.3797) ellipse (0.7761cm and 0.7761cm);



    \node[anchor=south] (text9533) at (6.7423, 18.2316){0};



    \path[draw=black] (2.9771, 19.9489).. controls (3.6971, 19.6462) and (4.8175, 19.1749) .. (5.6423, 18.8278);



    \path[draw=black,fill=black] (5.6829, 18.9445) -- (5.9602, 18.694) -- (5.5869, 18.717) -- (5.6829, 18.9445) -- (5.6829, 18.9445)-- cycle;;



    \node[anchor=south] (text9884) at (4.3716, 19.793){$\varepsilon$};



    \path[draw=black] (10.77, 24.8003) ellipse (0.8812cm and 0.8812cm);



    \path[draw=black] (10.77, 24.8003) ellipse (1.0224cm and 1.0224cm);



    \node[anchor=south] (text6845) at (10.77, 24.6521){$a_2$};



    \path[draw=black] (7.7354, 23.8375).. controls (8.236, 23.9984) and (8.8568, 24.1977) .. (9.4033, 24.3731);



    \path[draw=black,fill=black] (9.361, 24.4891) -- (9.7349, 24.4793) -- (9.4365, 24.2538) -- (9.361, 24.4891) -- (9.361, 24.4891)-- cycle;;



    \node[anchor=south] (text8743) at (8.7559, 24.427){a, b};



    \path[draw=black] (10.77, 22.2603) ellipse (0.8812cm and 0.8812cm);



    \node[anchor=south] (text4638) at (10.77, 22.1121){$a_3$};



    \path[draw=black] (7.7354, 23.223).. controls (8.2779, 23.0487) and (8.962, 22.8293) .. (9.5395, 22.6437);



    \path[draw=black,fill=black] (9.5702, 22.7637) -- (9.8683, 22.5383) -- (9.4947, 22.5284) -- (9.5702, 22.7637) -- (9.5702, 22.7637)-- cycle;;



    \node[anchor=south] (text1891) at (8.7559, 23.157){a};



    \path[draw=black] (10.4373, 25.7708).. controls (10.4225, 26.1458) and (10.5336, 26.4576) .. (10.77, 26.4576).. controls (10.9142, 26.4576) and (11.0116, 26.3419) .. (11.0624, 26.1666);



    \path[draw=black,fill=black] (11.1845, 26.187) -- (11.0973, 25.8237) -- (10.9386, 26.1624) -- (11.1845, 26.187) -- (11.1845, 26.187)-- cycle;;



    \node[anchor=south] (text5434) at (10.77, 26.6058){a, b};



    \path[draw=black] (10.77, 14.8872) ellipse (0.635cm and 0.635cm);



    \node[anchor=south] (text4540) at (10.77, 14.7391){1};



    \path[draw=black] (7.051, 17.6597).. controls (7.3127, 17.0596) and (7.7625, 16.2221) .. (8.3996, 15.7057).. controls (8.7852, 15.3928) and (9.2989, 15.1956) .. (9.7462, 15.0735);



    \path[draw=black,fill=black] (9.7663, 15.1956) -- (10.0806, 14.9934) -- (9.7088, 14.9553) -- (9.7663, 15.1956) -- (9.7663, 15.1956)-- cycle;;



    \node[anchor=south] (text9106) at (8.7559, 15.8538){a};



    \path[draw=black] (10.77, 20.1083) ellipse (0.635cm and 0.635cm);



    \node[anchor=south] (text7839) at (10.77, 19.9602){3};



    \path[draw=black] (7.4669, 18.681).. controls (8.1132, 18.9632) and (9.0851, 19.3876) .. (9.8009, 19.7005);



    \path[draw=black,fill=black] (9.7469, 19.8116) -- (10.1194, 19.8395) -- (9.8457, 19.5852) -- (9.7469, 19.8116) -- (9.7469, 19.8116)-- cycle;;



    \node[anchor=south] (text7823) at (8.7559, 19.5471){b};



    \path[draw=black] (25.8113, 20.3553) ellipse (0.635cm and 0.635cm);



    \node[anchor=south] (text2774) at (25.8113, 20.2071){8};



    \path[draw=black] (7.5212, 18.3797).. controls (8.3167, 18.3797) and (9.615, 18.3797) .. (10.7347, 18.3797).. controls (10.7347, 18.3797) and (10.7347, 18.3797) .. (22.9461, 18.3797).. controls (23.8157, 18.3797) and (24.6087, 19.0147) .. (25.1326, 19.5647);



    \path[draw=black,fill=black] (25.0317, 19.637) -- (25.3591, 19.8169) -- (25.2155, 19.4719) -- (25.0317, 19.637) -- (25.0317, 19.637)-- cycle;;



    \node[anchor=south] (text2999) at (17.2237, 18.5279){a};



    \path[draw=black] (22.9108, 20.3553) ellipse (0.635cm and 0.635cm);



    \path[draw=black] (22.9108, 20.3553) ellipse (0.7761cm and 0.7761cm);



    \node[anchor=south] (text9449) at (22.9108, 20.2071){7};



    \path[draw=black] (22.6328, 21.0958).. controls (22.5982, 21.4521) and (22.691, 21.7664) .. (22.9108, 21.7664).. controls (23.0413, 21.7664) and (23.1271, 21.6556) .. (23.168, 21.4909);



    \path[draw=black,fill=black] (23.2907, 21.5078) -- (23.1863, 21.149) -- (23.0442, 21.4948) -- (23.2907, 21.5078) -- (23.2907, 21.5078)-- cycle;;



    \node[anchor=south] (text5468) at (22.9108, 21.9146){a};



    \path[draw=black] (22.1238, 20.3122).. controls (21.5477, 20.283) and (20.7331, 20.2494) .. (20.018, 20.2494).. controls (14.1552, 20.2494) and (14.1552, 20.2494) .. (14.1552, 20.2494).. controls (13.3675, 20.2494) and (12.476, 20.2099) .. (11.811, 20.1729);



    \path[draw=black,fill=black] (11.823, 20.0498) -- (11.4635, 20.1524) -- (11.8082, 20.2964) -- (11.823, 20.0498) -- (11.823, 20.0498)-- cycle;;



    \node[anchor=south] (text5075) at (17.2237, 20.3976){b};



    \path[draw=black] (23.7046, 20.3553).. controls (24.028, 20.3553) and (24.4105, 20.3553) .. (24.7583, 20.3553);



    \path[draw=black,fill=black] (24.7523, 20.4788) -- (25.1051, 20.3553) -- (24.7523, 20.2318) -- (24.7523, 20.4788) -- (24.7523, 20.4788)-- cycle;;



    \node[anchor=south] (text268) at (24.4316, 20.5034){a};



    \path[draw=black,fill=black] (0.0635, 20.2494) ellipse (0.0635cm and 0.0635cm);



    \path[draw=black] (0.1351, 20.2494).. controls (0.2508, 20.2494) and (0.6262, 20.2494) .. (1.0255, 20.2494);



    \path[draw=black,fill=black] (1.0174, 20.3729) -- (1.3702, 20.2494) -- (1.0174, 20.126) -- (1.0174, 20.3729) -- (1.0174, 20.3729)-- cycle;;



    \path[draw=black] (14.1905, 25.6117) ellipse (0.8812cm and 0.8812cm);



    \node[anchor=south] (text9680) at (14.1905, 25.4635){$a_4$};



    \path[draw=black] (11.4198, 22.8748).. controls (11.9391, 23.3948) and (12.6898, 24.1455) .. (13.2733, 24.7294);



    \path[draw=black,fill=black] (13.1703, 24.801) -- (13.5068, 24.9633) -- (13.3449, 24.6264) -- (13.1703, 24.801) -- (13.1703, 24.801)-- cycle;;



    \node[anchor=south] (text2459) at (12.5508, 24.1847){b};



    \path[draw=black] (13.5488, 26.2396).. controls (13.106, 26.6506) and (12.4675, 27.1512) .. (11.7923, 27.3756).. controls (10.9298, 27.662) and (10.5632, 27.7763) .. (9.7476, 27.3756).. controls (8.6695, 26.846) and (7.8719, 25.7161) .. (7.3762, 24.8172);



    \path[draw=black,fill=black] (7.4874, 24.7629) -- (7.2136, 24.5085) -- (7.2686, 24.8782) -- (7.4874, 24.7629) -- (7.4874, 24.7629)-- cycle;;



    \node[anchor=south] (text1961) at (10.77, 27.783){a};



    \path[draw=black] (14.1905, 16.2278) ellipse (0.635cm and 0.635cm);



    \node[anchor=south] (text8565) at (14.1905, 16.0796){2};



    \path[draw=black] (11.3736, 15.1148).. controls (11.8787, 15.3169) and (12.6273, 15.6164) .. (13.2207, 15.8538);



    \path[draw=black,fill=black] (13.1611, 15.9628) -- (13.5347, 15.9794) -- (13.2528, 15.7339) -- (13.1611, 15.9628) -- (13.1611, 15.9628)-- cycle;;



    \node[anchor=south] (text7232) at (12.5508, 15.7459){b};



    \path[draw=black] (14.1905, 21.9781) ellipse (0.635cm and 0.635cm);



    \node[anchor=south] (text8946) at (14.1905, 21.8299){4};



    \path[draw=black] (11.3436, 20.4089).. controls (11.8611, 20.6978) and (12.6531, 21.1399) .. (13.2648, 21.481);



    \path[draw=black,fill=black] (13.2034, 21.5882) -- (13.5717, 21.6524) -- (13.3237, 21.3727) -- (13.2034, 21.5882) -- (13.2034, 21.5882)-- cycle;;



    \node[anchor=south] (text7844) at (12.5508, 21.2475){a};



    \path[draw=black] (17.2237, 22.1192) ellipse (0.635cm and 0.635cm);



    \node[anchor=south] (text4813) at (17.2237, 21.971){5};



    \path[draw=black] (25.466, 20.9038).. controls (25.115, 21.4641) and (24.487, 22.3044) .. (23.6869, 22.6836).. controls (21.6013, 23.6724) and (20.713, 23.3892) .. (18.4937, 22.7542).. controls (18.3748, 22.7199) and (18.2538, 22.6741) .. (18.1363, 22.6226);



    \path[draw=black,fill=black] (18.1995, 22.516) -- (17.8287, 22.4691) -- (18.0894, 22.7369) -- (18.1995, 22.516) -- (18.1995, 22.516)-- cycle;;



    \node[anchor=south] (text7149) at (21.3762, 23.4788){a};



    \path[draw=black] (13.5505, 16.3513).. controls (12.6958, 16.5308) and (11.0889, 16.8871) .. (9.7476, 17.2932).. controls (9.1073, 17.4872) and (8.4032, 17.7426) .. (7.8415, 17.9571);



    \path[draw=black,fill=black] (7.8123, 17.8361) -- (7.5279, 18.0785) -- (7.9015, 18.0665) -- (7.8123, 17.8361) -- (7.8123, 17.8361)-- cycle;;



    \node[anchor=south] (text6077) at (10.77, 17.4413){a};



    \path[draw=black] (14.8396, 22.0073).. controls (15.2248, 22.0257) and (15.7311, 22.0497) .. (16.1749, 22.0708);



    \path[draw=black,fill=black] (16.1636, 22.194) -- (16.522, 22.0874) -- (16.1756, 21.9474) -- (16.1636, 22.194) -- (16.1636, 22.194)-- cycle;;



    \node[anchor=south] (text1354) at (15.8302, 22.2052){b};



    \path[draw=black] (19.9827, 21.8722) ellipse (0.635cm and 0.635cm);



    \node[anchor=south] (text6769) at (19.9827, 21.7241){6};



    \path[draw=black] (17.8689, 22.0631).. controls (18.1853, 22.0341) and (18.5797, 21.9978) .. (18.9399, 21.9647);



    \path[draw=black,fill=black] (18.9466, 22.0881) -- (19.2867, 21.9329) -- (18.9241, 21.8422) -- (18.9466, 22.0881) -- (18.9466, 22.0881)-- cycle;;



    \node[anchor=south] (text6367) at (18.6034, 22.1513){a};



    \path[draw=black] (20.5557, 21.5868).. controls (20.9218, 21.3924) and (21.4182, 21.1289) .. (21.8581, 20.8954);



    \path[draw=black,fill=black] (21.9089, 21.0083) -- (22.1626, 20.7338) -- (21.7932, 20.7903) -- (21.9089, 21.0083) -- (21.9089, 21.0083)-- cycle;;



    \node[anchor=south] (text2484) at (21.3762, 21.3526){b};



  \end{scope}

\end{tikzpicture}

      \end{adjustbox}
      \caption{ПКА}
      \label{fig:afa}
    \end{figure}

    % \begin{figure}[H]
    %   \centering
    %   \begin{tabular}{|c|c|c|c|c|}
    %     \hline
    %     & bb & ab & a & bababa \\
    %     \hline
    %     aaabaa & $-$ & $+$ & $+$ & $+$ \\
    %     \hline
    %     aaab & $-$ & $-$ & $+$ & $+$ \\
    %     \hline
    %     $\varepsilon$ & $-$ & $-$ & $-$ & $+$ \\
    %     \hline
    %     aa & $-$ & $+$ & $-$ & $-$ \\
    %     \hline
    %   \end{tabular}
    %   \caption{Таблица префиксов-суффиксов}
    %   \label{table:afa-prefix-suffix}
    % \end{figure}

    \begin{figure}[H]
      \centering
      \begin{tabular}{|c|c|c|c|c|c|}
        \hline
        & bb & b & ab & $\varepsilon$ & a \\
        \hline
        aaabaaa & $-$ & $+$ & $+$ & $+$ & $+$ \\
        \hline
        aaabaa & $-$ & $-$ & $+$ & $+$ & $+$ \\
        \hline
        aaab & $-$ & $-$ & $-$ & $+$ & $+$ \\
        \hline
        ab & $-$ & $-$ & $-$ & $-$ & $+$\\
        \hline
        $\varepsilon$ & $-$ & $-$ & $-$ & $+$ & $-$\\
        \hline
      \end{tabular}
      \caption{Таблица префиксов-суффиксов}
      \label{table:afa-prefix-suffix}
    \end{figure}

    %% Получается, что в минимальном ПКА должно быть не менее 5? состояний, что не противоречит оценке выше.
    %% Таблицу получше "не придумаю" (код вряд ли сгенерирует, а вручную вообще не представляю как это 
    %% составлять).

    \section{Расширенное регулярное выражение}

    Преобразуем регулярное выражение:

    %% Здесь не то, чтобы можно чего изменить, поэтому просто что-то произвольное и бесполезное

    \[
      (aba)^{*}((aa | bab)aba^{*})^{*}
    \]

    Заметим, что любая непустая строка содержит букву $a$, а значит разделим регулярку на 2 альтернативы: пустая строка и непустая.

    \[
      (aba)^{*}((aa | bab)aba^{*})^{*})|\varepsilon
    \]

    Для непустого случая добавим проверку на наличие буквы $a$. Сделаем её через префиксную проверку $.^*a$.

    \[
      (?=.^*a)(aba)^{*}((aa | bab)aba^{*})^{*})|\varepsilon
    \]

    Заменив альтернативу с $\varepsilon$ и добавив маркеры начала и конца, получаем следующее расширенное регулярное выражение:

    \[
      \string^((?=.^*a)(aba)^{*}((aa | bab)aba^{*})^{*})?\$
    \]

    Следующий ПКА распознаёт указанное выше регулярное выражение:

    \begin{figure}[H]
      \centering
      \begin{adjustbox}{width=\textwidth}
      \def \globalscale {1.000000}
\begin{tikzpicture}[y=1cm, x=1cm, yscale=\globalscale,xscale=\globalscale, every node/.append style={scale=\globalscale}, inner sep=0pt, outer sep=0pt]
  \begin{scope}[shift={(0.1411, -11.2818)}]
    \path[fill=white] (-0.1411, 11.2889) -- (-0.1411, 22.7118) -- (24.6447, 22.7118) -- (24.6447, 11.2889) -- (-0.1411, 11.2889) -- (-0.1411, 11.2889)-- cycle;;



    \path[draw=black] (2.2278, 15.9456) ellipse (0.6544cm and 0.6544cm);



    \path[draw=black] (2.2278, 15.9456) ellipse (0.7955cm and 0.7955cm);



    \node[anchor=south] (text6836) at (2.2278, 15.7974){\&};



    \path[draw=black] (6.4964, 18.4856) ellipse (0.635cm and 0.635cm);



    \path[draw=black] (6.4964, 18.4856) ellipse (0.7761cm and 0.7761cm);



    \node[anchor=south] (text7508) at (6.4964, 18.3374){0};



    \path[draw=black] (2.9281, 16.3481).. controls (3.6139, 16.7629) and (4.6905, 17.4145) .. (5.4751, 17.889);



    \path[draw=black,fill=black] (5.4014, 17.9888) -- (5.7672, 18.0658) -- (5.5291, 17.7775) -- (5.4014, 17.9888) -- (5.4014, 17.9888)-- cycle;;



    \node[anchor=south] (text4701) at (4.3716, 17.7398){$\varepsilon$};



    \path[draw=black] (6.4964, 13.4408) ellipse (0.6354cm and 0.6354cm);



    \path[draw=black] (6.4964, 13.4408) ellipse (0.7765cm and 0.7765cm);



    \node[anchor=south] (text6637) at (6.4964, 13.2927){A};



    \path[draw=black] (2.9281, 15.5487).. controls (3.6107, 15.1416) and (4.681, 14.503) .. (5.4645, 14.0353);



    \path[draw=black,fill=black] (5.5273, 14.1418) -- (5.7669, 13.855) -- (5.4007, 13.9298) -- (5.5273, 14.1418) -- (5.5273, 14.1418)-- cycle;;



    \node[anchor=south] (text1049) at (4.3716, 15.2121){$\varepsilon$};



    \path[draw=black] (9.4248, 20.3553) ellipse (0.635cm and 0.635cm);



    \node[anchor=south] (text6908) at (9.4248, 20.2071){1};



    \path[draw=black] (7.166, 18.9004).. controls (7.5706, 19.1654) and (8.1005, 19.5118) .. (8.5411, 19.8004);



    \path[draw=black,fill=black] (8.4642, 19.8974) -- (8.8269, 19.9873) -- (8.5993, 19.6906) -- (8.4642, 19.8974) -- (8.4642, 19.8974)-- cycle;;



    \node[anchor=south] (text3296) at (8.0313, 19.6804){a};



    \path[draw=black] (9.4248, 16.5806) ellipse (0.635cm and 0.635cm);



    \node[anchor=south] (text1826) at (9.4248, 16.4324){3};



    \path[draw=black] (7.1049, 17.9984).. controls (7.3462, 17.8029) and (7.6341, 17.581) .. (7.9079, 17.399).. controls (8.0945, 17.2748) and (8.3012, 17.1524) .. (8.4998, 17.0409);



    \path[draw=black,fill=black] (8.5439, 17.1577) -- (8.7951, 16.8808) -- (8.4261, 16.9404) -- (8.5439, 17.1577) -- (8.5439, 17.1577)-- cycle;;



    \node[anchor=south] (text8626) at (8.0313, 17.5472){b};



    \path[draw=black] (23.8686, 16.8628) ellipse (0.635cm and 0.635cm);



    \node[anchor=south] (text5032) at (23.8686, 16.7146){8};



    \path[draw=black] (6.7649, 19.2264).. controls (7.1392, 20.2516) and (7.9791, 21.9781) .. (9.3895, 21.9781).. controls (9.3895, 21.9781) and (9.3895, 21.9781) .. (21.0033, 21.9781).. controls (22.8501, 21.9781) and (23.5013, 19.3936) .. (23.7232, 17.8989);



    \path[draw=black,fill=black] (23.8453, 17.9193) -- (23.7695, 17.5532) -- (23.6005, 17.8865) -- (23.8453, 17.9193) -- (23.8453, 17.9193)-- cycle;;



    \node[anchor=south] (text7443) at (15.2809, 22.1262){a};



    \path[draw=black] (20.9681, 16.8628) ellipse (0.635cm and 0.635cm);



    \path[draw=black] (20.9681, 16.8628) ellipse (0.7761cm and 0.7761cm);



    \node[anchor=south] (text4100) at (20.9681, 16.7146){7};



    \path[draw=black] (20.6901, 17.6033).. controls (20.6555, 17.9596) and (20.7483, 18.2739) .. (20.9681, 18.2739).. controls (21.0986, 18.2739) and (21.1843, 18.1631) .. (21.2252, 17.9984);



    \path[draw=black,fill=black] (21.348, 18.0153) -- (21.2436, 17.6565) -- (21.1014, 18.0023) -- (21.348, 18.0153) -- (21.348, 18.0153)-- cycle;;



    \node[anchor=south] (text9478) at (20.9681, 18.4221){a};



    \path[draw=black] (20.2297, 17.1514).. controls (19.6592, 17.3595) and (18.8288, 17.6036) .. (18.0753, 17.6036).. controls (12.3176, 17.6036) and (12.3176, 17.6036) .. (12.3176, 17.6036).. controls (11.624, 17.6036) and (10.8881, 17.333) .. (10.3371, 17.0702);



    \path[draw=black,fill=black] (10.4055, 16.9668) -- (10.0351, 16.9171) -- (10.2941, 17.187) -- (10.4055, 16.9668) -- (10.4055, 16.9668)-- cycle;;



    \node[anchor=south] (text4331) at (15.2809, 17.7518){b};



    \path[draw=black] (21.7618, 16.8628).. controls (22.0853, 16.8628) and (22.4677, 16.8628) .. (22.8152, 16.8628);



    \path[draw=black,fill=black] (22.8096, 16.9863) -- (23.1623, 16.8628) -- (22.8096, 16.7393) -- (22.8096, 16.9863) -- (22.8096, 16.9863)-- cycle;;



    \node[anchor=south] (text1255) at (22.4889, 17.0109){a};



    \path[draw=black] (12.3529, 12.2061) ellipse (0.635cm and 0.635cm);



    \path[draw=black] (12.3529, 12.2061) ellipse (0.7761cm and 0.7761cm);



    \node[anchor=south] (text6455) at (12.3529, 12.0579){B};



    \path[draw=black] (7.0838, 12.9208).. controls (7.5124, 12.555) and (8.1407, 12.0964) .. (8.7898, 11.8957).. controls (9.5754, 11.6526) and (10.5029, 11.7482) .. (11.2099, 11.895);



    \path[draw=black,fill=black] (11.1742, 12.0135) -- (11.5457, 11.9733) -- (11.23, 11.7729) -- (11.1742, 12.0135) -- (11.1742, 12.0135)-- cycle;;



    \node[anchor=south] (text4644) at (9.4248, 12.0438){a};



    \path[draw=black] (9.4248, 13.4408) ellipse (0.635cm and 0.635cm);



    \node[anchor=south] (text9180) at (9.4248, 13.2927){C};



    \path[draw=black] (7.2824, 13.4408).. controls (7.6175, 13.4408) and (8.0176, 13.4408) .. (8.3788, 13.4408);



    \path[draw=black,fill=black] (8.3721, 13.5643) -- (8.7249, 13.4408) -- (8.3721, 13.3174) -- (8.3721, 13.5643) -- (8.3721, 13.5643)-- cycle;;



    \node[anchor=south] (text1588) at (8.0313, 13.589){b};



    \path[draw=black] (12.0721, 12.9466).. controls (12.0375, 13.3029) and (12.131, 13.6172) .. (12.3529, 13.6172).. controls (12.4845, 13.6172) and (12.5709, 13.5065) .. (12.6122, 13.3417);



    \path[draw=black,fill=black] (12.7349, 13.359) -- (12.6305, 12.9999) -- (12.4883, 13.3456) -- (12.7349, 13.359) -- (12.7349, 13.359)-- cycle;;



    \node[anchor=south] (text2282) at (12.3529, 13.7654){a, b};



    \path[draw=black,fill=black] (0.0635, 15.9456) ellipse (0.0635cm and 0.0635cm);



    \path[draw=black] (0.1351, 15.9456).. controls (0.2508, 15.9456) and (0.6262, 15.9456) .. (1.0255, 15.9456);



    \path[draw=black,fill=black] (1.0174, 16.069) -- (1.3702, 15.9456) -- (1.0174, 15.8221) -- (1.0174, 16.069) -- (1.0174, 16.069)-- cycle;;



    \path[draw=black] (10.0245, 13.1971).. controls (10.3759, 13.045) and (10.8384, 12.8453) .. (11.2561, 12.6647);



    \path[draw=black,fill=black] (11.2988, 12.7808) -- (11.5739, 12.5275) -- (11.201, 12.554) -- (11.2988, 12.7808) -- (11.2988, 12.7808)-- cycle;;



    \node[anchor=south] (text6295) at (10.8183, 13.0457){a};



    \path[draw=black] (9.1779, 14.042).. controls (9.1274, 14.3884) and (9.2096, 14.7108) .. (9.4248, 14.7108).. controls (9.549, 14.7108) and (9.6291, 14.6029) .. (9.6647, 14.4452);



    \path[draw=black,fill=black] (9.7882, 14.4501) -- (9.6707, 14.0952) -- (9.5412, 14.4459) -- (9.7882, 14.4501) -- (9.7882, 14.4501)-- cycle;;



    \node[anchor=south] (text7571) at (9.4248, 14.859){b};



    \path[draw=black] (12.3529, 19.7908) ellipse (0.635cm and 0.635cm);



    \node[anchor=south] (text5402) at (12.3529, 19.6427){2};



    \path[draw=black] (10.0662, 20.2357).. controls (10.4299, 20.1637) and (10.9012, 20.0706) .. (11.3189, 19.988);



    \path[draw=black,fill=black] (11.3379, 20.1101) -- (11.66, 19.9207) -- (11.2899, 19.8681) -- (11.3379, 20.1101) -- (11.3379, 20.1101)-- cycle;;



    \node[anchor=south] (text6410) at (10.8183, 20.2551){b};



    \path[draw=black] (12.3529, 15.9103) ellipse (0.635cm and 0.635cm);



    \node[anchor=south] (text8306) at (12.3529, 15.7621){4};



    \path[draw=black] (10.052, 16.4416).. controls (10.4214, 16.3551) and (10.9054, 16.2415) .. (11.3316, 16.1417);



    \path[draw=black,fill=black] (11.3492, 16.2641) -- (11.6642, 16.0634) -- (11.2928, 16.0239) -- (11.3492, 16.2641) -- (11.3492, 16.2641)-- cycle;;



    \node[anchor=south] (text2839) at (10.8183, 16.4338){a};



    \path[draw=black] (15.2809, 15.6986) ellipse (0.635cm and 0.635cm);



    \node[anchor=south] (text8455) at (15.2809, 15.5504){5};



    \path[draw=black] (23.3176, 16.547).. controls (22.9045, 16.3057) and (22.3037, 15.9794) .. (21.7442, 15.7692).. controls (20.4308, 15.2753) and (20.0685, 15.1941) .. (18.675, 15.0283).. controls (18.1144, 14.9617) and (17.9627, 14.9412) .. (17.405, 15.0283).. controls (17.0169, 15.089) and (16.6024, 15.2097) .. (16.2458, 15.3321);



    \path[draw=black,fill=black] (16.2116, 15.2132) -- (15.9223, 15.4499) -- (16.2959, 15.4453) -- (16.2116, 15.2132) -- (16.2116, 15.2132)-- cycle;;



    \node[anchor=south] (text1998) at (19.4335, 15.2912){a};



    \path[draw=black] (11.7782, 19.4984).. controls (11.3312, 19.2705) and (10.6726, 18.9671) .. (10.0598, 18.8101).. controls (9.2851, 18.6115) and (8.3873, 18.5311) .. (7.6913, 18.4997);



    \path[draw=black,fill=black] (7.6969, 18.3765) -- (7.3399, 18.488) -- (7.6884, 18.6231) -- (7.6969, 18.3765) -- (7.6969, 18.3765)-- cycle;;



    \node[anchor=south] (text3681) at (9.4248, 18.9583){a};



    \path[draw=black] (12.9942, 15.8655).. controls (13.3555, 15.8387) and (13.8229, 15.8041) .. (14.2385, 15.773);



    \path[draw=black,fill=black] (14.2346, 15.8972) -- (14.5775, 15.748) -- (14.2166, 15.651) -- (14.2346, 15.8972) -- (14.2346, 15.8972)-- cycle;;



    \node[anchor=south] (text4001) at (13.8875, 15.9526){b};



    \path[draw=black] (18.04, 15.9808) ellipse (0.635cm and 0.635cm);



    \node[anchor=south] (text4984) at (18.04, 15.8327){6};



    \path[draw=black] (15.9262, 15.7628).. controls (16.2426, 15.796) and (16.637, 15.8373) .. (16.9968, 15.875);



    \path[draw=black,fill=black] (16.9803, 15.9974) -- (17.344, 15.9113) -- (17.006, 15.7519) -- (16.9803, 15.9974) -- (16.9803, 15.9974)-- cycle;;



    \node[anchor=south] (text4087) at (16.6606, 15.9964){a};



    \path[draw=black] (18.6535, 16.1593).. controls (18.9911, 16.2634) and (19.4282, 16.3985) .. (19.8293, 16.5223);



    \path[draw=black,fill=black] (19.7887, 16.6388) -- (20.1623, 16.625) -- (19.8617, 16.4031) -- (19.7887, 16.6388) -- (19.7887, 16.6388)-- cycle;;



    \node[anchor=south] (text7464) at (19.4335, 16.57){b};



  \end{scope}

\end{tikzpicture}

      \end{adjustbox}
      \caption{ПКА для расширенного регулярного выражения}
      \label{fig:afa-extended}
    \end{figure}

    Видим, что верхняя часть автомата распознаёт регулярное выражение $(aba)^{*}((aa | bab)aba^{*})^{*}$, а нижняя --- $.^* a .^*$, так как, 
    \[
      (?=.^*a)(aba)^{*}((aa | bab)aba^{*})^{*} = ((aba)^{*}((aa | bab)aba^{*})^{*}) \cap (.^* a .^*)
    \]
    Условие $\re?$ выполняется, так как состояния, в которые можно попасть после $\varepsilon$-перехода, являются финальными.

\end{document}

